\documentclass{article}
\usepackage[T1]{fontenc}
\usepackage{lmodern}
\usepackage{graphicx}
\usepackage{amsmath,amssymb}
\usepackage[a4paper,margin=2cm]{geometry}
\usepackage[absolute,overlay]{textpos}
\setlength{\TPHorizModule}{1mm}
\setlength{\TPVertModule}{1mm}
\usepackage[indonesian]{babel}
\usepackage{hyperref}
\setlength{\emergencystretch}{2em}
% unit: Halaman 26

\begin{document}

% === Halaman 26 (llm) ===

\section*{4. Forgery Attack}

\begin{itemize}
    \item Tujuan: membuat pesan palsu yang diterima sebagai pesan sah oleh sistem.
    \item Serangan ini sangat penting pada sistem yang menggunakan MAC (\textit{message authentication code}) atau \textit{digital signature} (akan dijelaskan pada materi kuliah selanjutnya).
    \item Misalkan Alice mengirim Pesan + MAC kepada Bob.
    \item Bob memeriksa MAC untuk memastikan pesan benar-benar berasal dari pihak yang memiliki kunci.
    \item Penyerang, Eve, ingin membuat:
    
    Pesan palsu + MAC yang valid
    
    sehingga Bob menganggap pesan tersebut autentik.
\end{itemize}

\end{document}
